\section{Second-Order}

In this section, we start with a concrete problem, then introduce second-order containers with interpretation, then the category of second-order containers, then second-order W-type to fix the problem.

\paragraph*{Bush Problem}

From indexed type to higher-order type:

How to understand the \AgdaDatatype{Bush} as example:

\begin{code}%
\>[0]\AgdaKeyword{data}\AgdaSpace{}%
\AgdaDatatype{Bush}\AgdaSpace{}%
\AgdaSymbol{(}\AgdaBound{X}\AgdaSpace{}%
\AgdaSymbol{:}\AgdaSpace{}%
\AgdaPrimitive{Set}\AgdaSymbol{)}\AgdaSpace{}%
\AgdaSymbol{:}\AgdaSpace{}%
\AgdaPrimitive{Set}\AgdaSpace{}%
\AgdaKeyword{where}\<%
\\
\>[0][@{}l@{\AgdaIndent{0}}]%
\>[2]\AgdaInductiveConstructor{nil}\AgdaSpace{}%
\AgdaSymbol{:}\AgdaSpace{}%
\AgdaDatatype{Bush}\AgdaSpace{}%
\AgdaBound{X}\<%
\\
%
\>[2]\AgdaInductiveConstructor{cons}\AgdaSpace{}%
\AgdaSymbol{:}\AgdaSpace{}%
\AgdaBound{X}\AgdaSpace{}%
\AgdaSymbol{→}\AgdaSpace{}%
\AgdaDatatype{Bush}\AgdaSpace{}%
\AgdaSymbol{(}\AgdaDatatype{Bush}\AgdaSpace{}%
\AgdaBound{X}\AgdaSymbol{)}\AgdaSpace{}%
\AgdaSymbol{→}\AgdaSpace{}%
\AgdaDatatype{Bush}\AgdaSpace{}%
\AgdaBound{X}\<%
\end{code}


If we naively encode this using indexed type:

\begin{code}
\>[0]\AgdaKeyword{data}\AgdaSpace{}%
\AgdaDatatype{Bush}\AgdaSpace{}%
\AgdaSymbol{:}\AgdaSpace{}%
\AgdaFunction{Type}\AgdaSpace{}%
\AgdaSymbol{→}\AgdaSpace{}%
\AgdaFunction{Type₁}\AgdaSpace{}%
\AgdaKeyword{where}\<%
\\
\>[0][@{}l@{\AgdaIndent{0}}]%
\>[2]\AgdaInductiveConstructor{nil}\AgdaSpace{}%
\AgdaSymbol{:}\AgdaSpace{}%
\AgdaSymbol{\{}\AgdaBound{X}\AgdaSpace{}%
\AgdaSymbol{:}\AgdaSpace{}%
\AgdaFunction{Type}\AgdaSymbol{\}}\AgdaSpace{}%
\AgdaSymbol{→}\AgdaSpace{}%
\AgdaDatatype{Bush}\AgdaSpace{}%
\AgdaBound{X}\<%
\\
%
\>[2]\AgdaInductiveConstructor{cons}\AgdaSpace{}%
\AgdaSymbol{:}\AgdaSpace{}%
\AgdaSymbol{\{}\AgdaBound{X}\AgdaSpace{}%
\AgdaSymbol{:}\AgdaSpace{}%
\AgdaFunction{Type}\AgdaSymbol{\}}\AgdaSpace{}%
\AgdaSymbol{→}\AgdaSpace{}%
\AgdaBound{X}\AgdaSpace{}%
\AgdaSymbol{→}\AgdaSpace{}%
\AgdaDatatype{Bush}\AgdaSpace{}%
\AgdaSymbol{(}\AgdaDatatype{Bush}\AgdaSpace{}%
\AgdaBound{X}\AgdaSymbol{)}\AgdaSpace{}%
\AgdaSymbol{→}\AgdaSpace{}%
\AgdaDatatype{Bush}\AgdaSpace{}%
\AgdaBound{X}\<%
\end{code}


The constructors \AgdaDatatype{nil} and \AgdaDatatype{cons} force it to be large type because they contains small universe \AgdaPrimitive{Type}. There are two issues here: there is a size issue in \AgdaDatatype{Bush (Bush X)}; there is a no positivity issue in \AgdaDatatype{Bush (Bush X)} because Agda cannot see it is algebraic.

\paragraph*{Second-Order Containers}

We need a new way to understand this. We introduce second-order containers

\begin{code}%
\>[0]\AgdaKeyword{record}\AgdaSpace{}%
\AgdaRecord{2Cont}\AgdaSpace{}%
\AgdaSymbol{:}\AgdaSpace{}%
\AgdaPrimitive{Set₁}\AgdaSpace{}%
\AgdaKeyword{where}\<%
\\
\>[0][@{}l@{\AgdaIndent{0}}]%
\>[2]\AgdaKeyword{pattern}\<%
\\
%
\>[2]\AgdaKeyword{inductive}\<%
\\
%
\>[2]\AgdaKeyword{constructor}\AgdaSpace{}%
\AgdaOperator{\AgdaInductiveConstructor{\AgdaUnderscore{}◃\AgdaUnderscore{}+\AgdaUnderscore{}+\AgdaUnderscore{}}}\<%
\\
%
\>[2]\AgdaKeyword{field}\<%
\\
\>[2][@{}l@{\AgdaIndent{0}}]%
\>[4]\AgdaField{S}\AgdaSpace{}%
\AgdaSymbol{:}\AgdaSpace{}%
\AgdaPrimitive{Set}\<%
\\
%
\>[4]\AgdaField{PX}\AgdaSpace{}%
\AgdaSymbol{:}\AgdaSpace{}%
\AgdaField{S}\AgdaSpace{}%
\AgdaSymbol{→}\AgdaSpace{}%
\AgdaPrimitive{Set}\<%
\\
%
\>[4]\AgdaField{PF}\AgdaSpace{}%
\AgdaSymbol{:}\AgdaSpace{}%
\AgdaField{S}\AgdaSpace{}%
\AgdaSymbol{→}\AgdaSpace{}%
\AgdaPrimitive{Set}\<%
\\
%
\>[4]\AgdaField{RF}\AgdaSpace{}%
\AgdaSymbol{:}\AgdaSpace{}%
\AgdaSymbol{(}\AgdaBound{s}\AgdaSpace{}%
\AgdaSymbol{:}\AgdaSpace{}%
\AgdaField{S}\AgdaSymbol{)}\AgdaSpace{}%
\AgdaSymbol{→}\AgdaSpace{}%
\AgdaField{PF}\AgdaSpace{}%
\AgdaBound{s}\AgdaSpace{}%
\AgdaSymbol{→}\AgdaSpace{}%
\AgdaRecord{2Cont}\<%
\end{code}


There is a function interpretation:

\begin{code}
\>[0]\AgdaOperator{\AgdaFunction{2⟦}}\AgdaSpace{}%
\AgdaBound{S}\AgdaSpace{}%
\AgdaOperator{\AgdaInductiveConstructor{◃}}\AgdaSpace{}%
\AgdaBound{PX}\AgdaSpace{}%
\AgdaOperator{\AgdaInductiveConstructor{+}}\AgdaSpace{}%
\AgdaBound{PF}\AgdaSpace{}%
\AgdaOperator{\AgdaInductiveConstructor{+}}\AgdaSpace{}%
\AgdaBound{RF}\AgdaSpace{}%
\AgdaOperator{\AgdaFunction{⟧}}\AgdaSpace{}%
\AgdaBound{F}\AgdaSpace{}%
\AgdaBound{X}\AgdaSpace{}%
\AgdaSymbol{=}\<%
\\
\>[0][@{}l@{\AgdaIndent{0}}]%
\>[2]\AgdaFunction{Σ[}\AgdaSpace{}%
\AgdaBound{s}\AgdaSpace{}%
\AgdaFunction{∈}\AgdaSpace{}%
\AgdaBound{S}\AgdaSpace{}%
\AgdaFunction{]}\AgdaSpace{}%
\AgdaSymbol{((}\AgdaBound{PX}\AgdaSpace{}%
\AgdaBound{s}\AgdaSpace{}%
\AgdaSymbol{→}\AgdaSpace{}%
\AgdaBound{X}\AgdaSymbol{)}\AgdaSpace{}%
\AgdaOperator{\AgdaFunction{×}}\AgdaSpace{}%
\AgdaSymbol{((}\AgdaBound{pF}\AgdaSpace{}%
\AgdaSymbol{:}\AgdaSpace{}%
\AgdaBound{PF}\AgdaSpace{}%
\AgdaBound{s}\AgdaSymbol{)}\AgdaSpace{}%
\AgdaSymbol{→}\AgdaSpace{}%
\AgdaBound{F}\AgdaSpace{}%
\AgdaSymbol{(}\AgdaOperator{\AgdaFunction{2⟦}}\AgdaSpace{}%
\AgdaBound{RF}\AgdaSpace{}%
\AgdaBound{s}\AgdaSpace{}%
\AgdaBound{pF}\AgdaSpace{}%
\AgdaOperator{\AgdaFunction{⟧}}\AgdaSpace{}%
\AgdaBound{F}\AgdaSpace{}%
\AgdaBound{X}\AgdaSymbol{)))}\<%
\end{code}



There is a categorical interpretation

\begin{code}
\>[0]\AgdaFunction{appS}\AgdaSpace{}%
\AgdaSymbol{:}\AgdaSpace{}%
\AgdaRecord{2Cont}\AgdaSpace{}%
\AgdaSymbol{→}\AgdaSpace{}%
\AgdaRecord{Cont}\AgdaSpace{}%
\AgdaSymbol{→}\AgdaSpace{}%
\AgdaPrimitive{Set}\<%
\\
\>[0]\AgdaFunction{appS}\AgdaSpace{}%
\AgdaSymbol{(}\AgdaBound{S}\AgdaSpace{}%
\AgdaOperator{\AgdaInductiveConstructor{◃}}\AgdaSpace{}%
\AgdaBound{PX}\AgdaSpace{}%
\AgdaOperator{\AgdaInductiveConstructor{+}}\AgdaSpace{}%
\AgdaBound{PF}\AgdaSpace{}%
\AgdaOperator{\AgdaInductiveConstructor{+}}\AgdaSpace{}%
\AgdaBound{RF}\AgdaSymbol{)}\AgdaSpace{}%
\AgdaSymbol{(}\AgdaBound{T}\AgdaSpace{}%
\AgdaOperator{\AgdaInductiveConstructor{◃}}\AgdaSpace{}%
\AgdaBound{Q}\AgdaSymbol{)}\AgdaSpace{}%
\AgdaSymbol{=}\<%
\\
\>[0][@{}l@{\AgdaIndent{0}}]%
\>[2]\AgdaFunction{Σ[}\AgdaSpace{}%
\AgdaBound{s}\AgdaSpace{}%
\AgdaFunction{∈}\AgdaSpace{}%
\AgdaBound{S}\AgdaSpace{}%
\AgdaFunction{]}\AgdaSpace{}%
\AgdaSymbol{((}\AgdaBound{pF}\AgdaSpace{}%
\AgdaSymbol{:}\AgdaSpace{}%
\AgdaBound{PF}\AgdaSpace{}%
\AgdaBound{s}\AgdaSymbol{)}\AgdaSpace{}%
\AgdaSymbol{→}\AgdaSpace{}%
\AgdaFunction{Σ[}\AgdaSpace{}%
\AgdaBound{t}\AgdaSpace{}%
\AgdaFunction{∈}\AgdaSpace{}%
\AgdaBound{T}\AgdaSpace{}%
\AgdaFunction{]}\AgdaSpace{}%
\AgdaSymbol{(}\AgdaBound{Q}\AgdaSpace{}%
\AgdaBound{t}\AgdaSpace{}%
\AgdaSymbol{→}\AgdaSpace{}%
\AgdaFunction{appS}\AgdaSpace{}%
\AgdaSymbol{(}\AgdaBound{RF}\AgdaSpace{}%
\AgdaBound{s}\AgdaSpace{}%
\AgdaBound{pF}\AgdaSymbol{)}\AgdaSpace{}%
\AgdaSymbol{(}\AgdaBound{T}\AgdaSpace{}%
\AgdaOperator{\AgdaInductiveConstructor{◃}}\AgdaSpace{}%
\AgdaBound{Q}\AgdaSymbol{)))}\<%
\\
%
\\[\AgdaEmptyExtraSkip]%
\>[0]\AgdaFunction{appP}\AgdaSpace{}%
\AgdaSymbol{:}\AgdaSpace{}%
\AgdaSymbol{(}\AgdaBound{H}\AgdaSpace{}%
\AgdaSymbol{:}\AgdaSpace{}%
\AgdaRecord{2Cont}\AgdaSymbol{)}\AgdaSpace{}%
\AgdaSymbol{(}\AgdaBound{F}\AgdaSpace{}%
\AgdaSymbol{:}\AgdaSpace{}%
\AgdaRecord{Cont}\AgdaSymbol{)}\AgdaSpace{}%
\AgdaSymbol{→}\AgdaSpace{}%
\AgdaFunction{appS}\AgdaSpace{}%
\AgdaBound{H}\AgdaSpace{}%
\AgdaBound{F}\AgdaSpace{}%
\AgdaSymbol{→}\AgdaSpace{}%
\AgdaPrimitive{Set}\<%
\\
\>[0]\AgdaFunction{appP}\AgdaSpace{}%
\AgdaSymbol{(}\AgdaBound{S}\AgdaSpace{}%
\AgdaOperator{\AgdaInductiveConstructor{◃}}\AgdaSpace{}%
\AgdaBound{PX}\AgdaSpace{}%
\AgdaOperator{\AgdaInductiveConstructor{+}}\AgdaSpace{}%
\AgdaBound{PF}\AgdaSpace{}%
\AgdaOperator{\AgdaInductiveConstructor{+}}\AgdaSpace{}%
\AgdaBound{RF}\AgdaSymbol{)}\AgdaSpace{}%
\AgdaSymbol{(}\AgdaBound{T}\AgdaSpace{}%
\AgdaOperator{\AgdaInductiveConstructor{◃}}\AgdaSpace{}%
\AgdaBound{Q}\AgdaSymbol{)}\AgdaSpace{}%
\AgdaSymbol{(}\AgdaBound{s}\AgdaSpace{}%
\AgdaOperator{\AgdaInductiveConstructor{,}}\AgdaSpace{}%
\AgdaBound{f}\AgdaSymbol{)}\AgdaSpace{}%
\AgdaSymbol{=}\<%
\\
\>[0][@{}l@{\AgdaIndent{0}}]%
\>[2]\AgdaBound{PX}\AgdaSpace{}%
\AgdaBound{s}\AgdaSpace{}%
\AgdaOperator{\AgdaDatatype{⊎}}\AgdaSpace{}%
\AgdaFunction{Σ[}\AgdaSpace{}%
\AgdaBound{pF}\AgdaSpace{}%
\AgdaFunction{∈}\AgdaSpace{}%
\AgdaBound{PF}\AgdaSpace{}%
\AgdaBound{s}\AgdaSpace{}%
\AgdaFunction{]}\AgdaSpace{}%
\AgdaKeyword{let}\AgdaSpace{}%
\AgdaSymbol{(}\AgdaBound{t}\AgdaSpace{}%
\AgdaOperator{\AgdaInductiveConstructor{,}}\AgdaSpace{}%
\AgdaBound{g}\AgdaSymbol{)}\AgdaSpace{}%
\AgdaSymbol{=}\AgdaSpace{}%
\AgdaBound{f}\AgdaSpace{}%
\AgdaBound{pF}\AgdaSpace{}%
\AgdaKeyword{in}\<%
\\
\>[2][@{}l@{\AgdaIndent{0}}]%
\>[4]\AgdaFunction{Σ[}\AgdaSpace{}%
\AgdaBound{q}\AgdaSpace{}%
\AgdaFunction{∈}\AgdaSpace{}%
\AgdaBound{Q}\AgdaSpace{}%
\AgdaBound{t}\AgdaSpace{}%
\AgdaFunction{]}\AgdaSpace{}%
\AgdaFunction{appP}\AgdaSpace{}%
\AgdaSymbol{(}\AgdaBound{RF}\AgdaSpace{}%
\AgdaBound{s}\AgdaSpace{}%
\AgdaBound{pF}\AgdaSymbol{)}\AgdaSpace{}%
\AgdaSymbol{(}\AgdaBound{T}\AgdaSpace{}%
\AgdaOperator{\AgdaInductiveConstructor{◃}}\AgdaSpace{}%
\AgdaBound{Q}\AgdaSymbol{)}\AgdaSpace{}%
\AgdaSymbol{(}\AgdaBound{g}\AgdaSpace{}%
\AgdaBound{q}\AgdaSymbol{)}\<%
\\
%
\\[\AgdaEmptyExtraSkip]%
\>[0]\AgdaFunction{app}\AgdaSpace{}%
\AgdaSymbol{:}\AgdaSpace{}%
\AgdaRecord{2Cont}\AgdaSpace{}%
\AgdaSymbol{→}\AgdaSpace{}%
\AgdaRecord{Cont}\AgdaSpace{}%
\AgdaSymbol{→}\AgdaSpace{}%
\AgdaRecord{Cont}\<%
\\
\>[0]\AgdaFunction{app}\AgdaSpace{}%
\AgdaBound{H}\AgdaSpace{}%
\AgdaBound{F}\AgdaSpace{}%
\AgdaSymbol{=}\AgdaSpace{}%
\AgdaFunction{appS}\AgdaSpace{}%
\AgdaBound{H}\AgdaSpace{}%
\AgdaBound{F}\AgdaSpace{}%
\AgdaOperator{\AgdaInductiveConstructor{◃}}\AgdaSpace{}%
\AgdaFunction{appP}\AgdaSpace{}%
\AgdaBound{H}\AgdaSpace{}%
\AgdaBound{F}\<%
\end{code}


The definition is derived by solving the following equation:

\[ 2⟦ H ⟧ ⟦ F ⟧ X ≃ ⟦ app H F ⟧ X \]

\AgdaFunction{2app}\AgdaSpace{}\AgdaBound{H} is an endofunctor of category of \AgdaRecord{Cont}. We omit the functorial action \AgdaFunction{app₁}.

\paragraph*{The category of second-order containers}

We present the morphisms of second-order containers:

\begin{code}
Some code here
\end{code}

They form a wild category. It has all infinitary products and coproducts, constructed pointwisely on shape.

\paragraph*{Second-Order W-type}

We want to calculate the least algebras of second-order container functors. That is we want to have ... :

\[
\begin{tikzcd}[column sep=large, row sep=large]
  \mathsf{app}\,\mathsf{H}\,(\mathsf{2W}\,\mathsf{H})
    \arrow[r, "\mathsf{app₁}\,\mathsf{H}\,(\mathsf{2fold}\,\alpha)"]
    \arrow[d, "{2sup}"'] &
  \mathsf{app₁}\,\mathsf{H}\,F
    \arrow[d, "\alpha"] \\
  \mathsf{2W}\,\mathsf{H}
    \arrow[r, dashed, "\mathsf{2fold}\,\alpha"'] &
  F
\end{tikzcd}
\]

It follows:

\[ 2W H ≃ app H (2W H) \]

We actually need to generalised the equation to:

\[ 2W-d H' H ≃ app H' (2W-d H' H) \]

because ...

We construct 2W by inductive-induction. It is constructed by splitting the domain \AgdaRecord{Cont} into \AgdaPrimitive{S : Set} and \AgdaPrimitive{P : S → Set}. 

\begin{code}%
\>[0]\AgdaKeyword{record}\AgdaSpace{}%
\AgdaRecord{2WS-d}\AgdaSpace{}%
\AgdaSymbol{(}\AgdaBound{H'}\AgdaSpace{}%
\AgdaBound{H}\AgdaSpace{}%
\AgdaSymbol{:}\AgdaSpace{}%
\AgdaRecord{2Cont}\AgdaSymbol{)}\AgdaSpace{}%
\AgdaSymbol{:}\AgdaSpace{}%
\AgdaPrimitive{Set}\<%
\\
%
\\[\AgdaEmptyExtraSkip]%
\>[0]\AgdaKeyword{data}\AgdaSpace{}%
\AgdaDatatype{2WP-d}\AgdaSpace{}%
\AgdaSymbol{(}\AgdaBound{H'}\AgdaSpace{}%
\AgdaBound{H}\AgdaSpace{}%
\AgdaSymbol{:}\AgdaSpace{}%
\AgdaRecord{2Cont}\AgdaSymbol{)}\AgdaSpace{}%
\AgdaSymbol{(}\AgdaBound{2ws}\AgdaSpace{}%
\AgdaSymbol{:}\AgdaSpace{}%
\AgdaRecord{2WS-d}\AgdaSpace{}%
\AgdaBound{H'}\AgdaSpace{}%
\AgdaBound{H}\AgdaSymbol{)}\AgdaSpace{}%
\AgdaSymbol{:}\AgdaSpace{}%
\AgdaPrimitive{Set}\<%
\\
%
\\[\AgdaEmptyExtraSkip]%
\>[0]\AgdaKeyword{record}\AgdaSpace{}%
\AgdaRecord{2WS-d}\AgdaSpace{}%
\AgdaBound{H'}\AgdaSpace{}%
\AgdaBound{H}\AgdaSpace{}%
\AgdaKeyword{where}\<%
\\
\>[0][@{}l@{\AgdaIndent{0}}]%
\>[2]\AgdaKeyword{constructor}\AgdaSpace{}%
\AgdaInductiveConstructor{mk}\<%
\\
%
\>[2]\AgdaKeyword{inductive}\<%
\\
%
\>[2]\AgdaKeyword{pattern}\<%
\\
%
\>[2]\AgdaKeyword{open}\AgdaSpace{}%
\AgdaModule{2Cont}\AgdaSpace{}%
\AgdaBound{H'}\<%
\\
%
\>[2]\AgdaKeyword{field}\<%
\\
\>[2][@{}l@{\AgdaIndent{0}}]%
\>[4]\AgdaField{s}\AgdaSpace{}%
\AgdaSymbol{:}\AgdaSpace{}%
\AgdaFunction{S}\<%
\\
%
\>[4]\AgdaField{f}\AgdaSpace{}%
\AgdaSymbol{:}\AgdaSpace{}%
\AgdaSymbol{(}\AgdaBound{pF}\AgdaSpace{}%
\AgdaSymbol{:}\AgdaSpace{}%
\AgdaFunction{PF}\AgdaSpace{}%
\AgdaField{s}\AgdaSymbol{)}\AgdaSpace{}%
\AgdaSymbol{→}\AgdaSpace{}%
\AgdaFunction{Σ[}\AgdaSpace{}%
\AgdaBound{t}\AgdaSpace{}%
\AgdaFunction{∈}\AgdaSpace{}%
\AgdaRecord{2WS-d}\AgdaSpace{}%
\AgdaBound{H}\AgdaSpace{}%
\AgdaBound{H}\AgdaSpace{}%
\AgdaFunction{]}\AgdaSpace{}%
\AgdaSymbol{(}\AgdaDatatype{2WP-d}\AgdaSpace{}%
\AgdaBound{H}\AgdaSpace{}%
\AgdaBound{H}\AgdaSpace{}%
\AgdaBound{t}\AgdaSpace{}%
\AgdaSymbol{→}\AgdaSpace{}%
\AgdaRecord{2WS-d}\AgdaSpace{}%
\AgdaSymbol{(}\AgdaFunction{RF}\AgdaSpace{}%
\AgdaField{s}\AgdaSpace{}%
\AgdaBound{pF}\AgdaSymbol{)}\AgdaSpace{}%
\AgdaBound{H}\AgdaSymbol{)}\<%
\\
%
\\[\AgdaEmptyExtraSkip]%
\>[0]\AgdaKeyword{data}\AgdaSpace{}%
\AgdaDatatype{2WP-d}\AgdaSpace{}%
\AgdaBound{H'}\AgdaSpace{}%
\AgdaBound{H}\AgdaSpace{}%
\AgdaBound{ws}\AgdaSpace{}%
\AgdaKeyword{where}\<%
\\
\>[0][@{}l@{\AgdaIndent{0}}]%
\>[2]\AgdaInductiveConstructor{posX}\AgdaSpace{}%
\AgdaSymbol{:}%
\>[67I]\AgdaSymbol{(}\AgdaKeyword{open}\AgdaSpace{}%
\AgdaModule{2Cont}\AgdaSpace{}%
\AgdaBound{H'}\AgdaSymbol{)}\AgdaSpace{}%
\AgdaSymbol{(}\AgdaKeyword{open}\AgdaSpace{}%
\AgdaModule{2WS-d}\AgdaSpace{}%
\AgdaBound{ws}\AgdaSymbol{)}\<%
\\
\>[.][@{}l@{}]\<[67I]%
\>[9]\AgdaSymbol{→}\AgdaSpace{}%
\AgdaFunction{PX}\AgdaSpace{}%
\AgdaField{s}\<%
\\
%
\>[9]\AgdaSymbol{→}\AgdaSpace{}%
\AgdaDatatype{2WP-d}\AgdaSpace{}%
\AgdaBound{H'}\AgdaSpace{}%
\AgdaBound{H}\AgdaSpace{}%
\AgdaBound{ws}\<%
\\
%
\>[2]\AgdaInductiveConstructor{posF}\AgdaSpace{}%
\AgdaSymbol{:}%
\>[80I]\AgdaSymbol{(}\AgdaKeyword{open}\AgdaSpace{}%
\AgdaModule{2Cont}\AgdaSpace{}%
\AgdaBound{H'}\AgdaSymbol{)}\AgdaSpace{}%
\AgdaSymbol{(}\AgdaKeyword{open}\AgdaSpace{}%
\AgdaModule{2WS-d}\AgdaSpace{}%
\AgdaBound{ws}\AgdaSymbol{)}\<%
\\
\>[.][@{}l@{}]\<[80I]%
\>[9]\AgdaSymbol{→}%
\>[86I]\AgdaSymbol{(}\AgdaFunction{Σ[}\AgdaSpace{}%
\AgdaBound{pF}\AgdaSpace{}%
\AgdaFunction{∈}\AgdaSpace{}%
\AgdaFunction{PF}\AgdaSpace{}%
\AgdaField{s}\AgdaSpace{}%
\AgdaFunction{]}\AgdaSpace{}%
\AgdaKeyword{let}\AgdaSpace{}%
\AgdaSymbol{(}\AgdaBound{t}\AgdaSpace{}%
\AgdaOperator{\AgdaInductiveConstructor{,}}\AgdaSpace{}%
\AgdaBound{g}\AgdaSymbol{)}\AgdaSpace{}%
\AgdaSymbol{=}\AgdaSpace{}%
\AgdaField{f}\AgdaSpace{}%
\AgdaBound{pF}\AgdaSpace{}%
\AgdaKeyword{in}\<%
\\
\>[86I][@{}l@{\AgdaIndent{0}}]%
\>[13]\AgdaFunction{Σ[}\AgdaSpace{}%
\AgdaBound{q}\AgdaSpace{}%
\AgdaFunction{∈}\AgdaSpace{}%
\AgdaDatatype{2WP-d}\AgdaSpace{}%
\AgdaBound{H}\AgdaSpace{}%
\AgdaBound{H}\AgdaSpace{}%
\AgdaBound{t}\AgdaSpace{}%
\AgdaFunction{]}\AgdaSpace{}%
\AgdaDatatype{2WP-d}\AgdaSpace{}%
\AgdaSymbol{(}\AgdaFunction{RF}\AgdaSpace{}%
\AgdaField{s}\AgdaSpace{}%
\AgdaBound{pF}\AgdaSymbol{)}\AgdaSpace{}%
\AgdaBound{H}\AgdaSpace{}%
\AgdaSymbol{(}\AgdaBound{g}\AgdaSpace{}%
\AgdaBound{q}\AgdaSymbol{))}\<%
\\
%
\>[9]\AgdaSymbol{→}\AgdaSpace{}%
\AgdaDatatype{2WP-d}\AgdaSpace{}%
\AgdaBound{H'}\AgdaSpace{}%
\AgdaBound{H}\AgdaSpace{}%
\AgdaBound{ws}\<%
\\
%
\\[\AgdaEmptyExtraSkip]%
\>[0]\AgdaFunction{2WS}\AgdaSpace{}%
\AgdaSymbol{:}\AgdaSpace{}%
\AgdaRecord{2Cont}\AgdaSpace{}%
\AgdaSymbol{→}\AgdaSpace{}%
\AgdaPrimitive{Set}\<%
\\
\>[0]\AgdaFunction{2WS}\AgdaSpace{}%
\AgdaBound{H}\AgdaSpace{}%
\AgdaSymbol{=}\AgdaSpace{}%
\AgdaRecord{2WS-d}\AgdaSpace{}%
\AgdaBound{H}\AgdaSpace{}%
\AgdaBound{H}\<%
\\
%
\\[\AgdaEmptyExtraSkip]%
\>[0]\AgdaFunction{2WP}\AgdaSpace{}%
\AgdaSymbol{:}\AgdaSpace{}%
\AgdaSymbol{(}\AgdaBound{H}\AgdaSpace{}%
\AgdaSymbol{:}\AgdaSpace{}%
\AgdaRecord{2Cont}\AgdaSymbol{)}\AgdaSpace{}%
\AgdaSymbol{→}\AgdaSpace{}%
\AgdaFunction{2WS}\AgdaSpace{}%
\AgdaBound{H}\AgdaSpace{}%
\AgdaSymbol{→}\AgdaSpace{}%
\AgdaPrimitive{Set}\<%
\\
\>[0]\AgdaFunction{2WP}\AgdaSpace{}%
\AgdaBound{H}\AgdaSpace{}%
\AgdaSymbol{=}\AgdaSpace{}%
\AgdaDatatype{2WP-d}\AgdaSpace{}%
\AgdaBound{H}\AgdaSpace{}%
\AgdaBound{H}\<%
\\
%
\\[\AgdaEmptyExtraSkip]%
\>[0]\AgdaFunction{2W-d}\AgdaSpace{}%
\AgdaSymbol{:}\AgdaSpace{}%
\AgdaRecord{2Cont}\AgdaSpace{}%
\AgdaSymbol{→}\AgdaSpace{}%
\AgdaRecord{2Cont}\AgdaSpace{}%
\AgdaSymbol{→}\AgdaSpace{}%
\AgdaRecord{Cont}\<%
\\
\>[0]\AgdaFunction{2W-d}\AgdaSpace{}%
\AgdaBound{H'}\AgdaSpace{}%
\AgdaBound{H}\AgdaSpace{}%
\AgdaSymbol{=}\AgdaSpace{}%
\AgdaRecord{2WS-d}\AgdaSpace{}%
\AgdaBound{H'}\AgdaSpace{}%
\AgdaBound{H}\AgdaSpace{}%
\AgdaOperator{\AgdaInductiveConstructor{◃}}\AgdaSpace{}%
\AgdaDatatype{2WP-d}\AgdaSpace{}%
\AgdaBound{H'}\AgdaSpace{}%
\AgdaBound{H}\<%
\\
%
\\[\AgdaEmptyExtraSkip]%
\>[0]\AgdaFunction{2W}\AgdaSpace{}%
\AgdaSymbol{:}\AgdaSpace{}%
\AgdaRecord{2Cont}\AgdaSpace{}%
\AgdaSymbol{→}\AgdaSpace{}%
\AgdaRecord{Cont}\<%
\\
\>[0]\AgdaFunction{2W}\AgdaSpace{}%
\AgdaBound{H}\AgdaSpace{}%
\AgdaSymbol{=}\AgdaSpace{}%
\AgdaFunction{2W-d}\AgdaSpace{}%
\AgdaBound{H}\AgdaSpace{}%
\AgdaBound{H}\<%
\end{code}


Then we can define some morphisms from the diagram:

\begin{code}
\>[2]\AgdaFunction{2sup}\AgdaSpace{}%
\AgdaSymbol{:}\AgdaSpace{}%
\AgdaFunction{app}\AgdaSpace{}%
\AgdaBound{H}\AgdaSpace{}%
\AgdaSymbol{(}\AgdaFunction{2W}\AgdaSpace{}%
\AgdaBound{H}\AgdaSymbol{)}\AgdaSpace{}%
\AgdaOperator{\AgdaRecord{→ᶜ}}\AgdaSpace{}%
\AgdaFunction{2W}\AgdaSpace{}%
\AgdaBound{H}\<%
\\
%
\>[2]\AgdaFunction{2sup}\AgdaSpace{}%
\AgdaSymbol{=}\AgdaSpace{}%
\AgdaFunction{2sup-d}\<%
\\
%
\\[\AgdaEmptyExtraSkip]%
%
\>[2]\AgdaFunction{fold2W}\AgdaSpace{}%
\AgdaSymbol{:}\AgdaSpace{}%
\AgdaFunction{2W}\AgdaSpace{}%
\AgdaBound{H}\AgdaSpace{}%
\AgdaOperator{\AgdaRecord{→ᶜ}}\AgdaSpace{}%
\AgdaBound{TQ}\<%
\\
%
\>[2]\AgdaFunction{fold2W}\AgdaSpace{}%
\AgdaSymbol{=}\AgdaSpace{}%
\AgdaFunction{fold2WS}\AgdaSpace{}%
\AgdaOperator{\AgdaInductiveConstructor{◃}}\AgdaSpace{}%
\AgdaFunction{fold2WP}\<%
\end{code}


We can show that digram commutes and \AgdaFunction{fold2W} is unique.

\paragraph*{Bush as 2W-type}

Now we are ready to construct \AgdaDatatype{Bush}. We define \AgdaFunction{Hᶜ : 2Cont}, where its interpretation is \AgdaFunction{2⟦ Hᶜ ⟧ F X ≃ ⊤ ⊎ X × F (F X)}. The container representation \AgdaFunction{Bushᶜ} can be derived by applying \AgdaFunction{2W} on \AgdaFunction{Hᶜ}. Finally, the original \AgdaDatatype{Bush} is just the extension functor \AgdaFunction{⟦ Bushᶜ ⟧}.
